\section{Proceso de muestreo de DDIM}

\subsection{Fórmula general de muestreo}

Dado una muestra ruidosa $\mathbf{x}_t$ en el paso temporal $t$, la fórmula de actualización por paso de DDIM es:

$$
\mathbf{x}_{t-1} = \sqrt{\bar{\alpha}_{t-1}}\,\underbrace{\hat{\mathbf{x}}_0(\mathbf{x}_t)}_{\text{datos limpios predichos}} + \sqrt{1-\bar{\alpha}_{t-1}-\sigma_t^2}\cdot\underbrace{\frac{\mathbf{x}_t - \sqrt{\bar{\alpha}_t}\,\hat{\mathbf{x}}_0(\mathbf{x}_t)}{\sqrt{1-\bar{\alpha}_t}}}_{\text{"dirección" hacia } \mathbf{x}_t} + \underbrace{\sigma_t\,\boldsymbol{\epsilon}_t}_{\text{ruido aleatorio}}
$$

donde $\hat{\mathbf{x}}_0(\mathbf{x}_t)$ es la estimación de $\mathbf{x}_0$ mediante una red de predicción de ruido:
$$
\hat{\mathbf{x}}_0(\mathbf{x}_t) = \frac{\mathbf{x}_t - \sqrt{1-\bar{\alpha}_t}\,\boldsymbol{\epsilon}_\theta(\mathbf{x}_t, t)}{\sqrt{\bar{\alpha}_t}}
$$

$\boldsymbol{\epsilon}_t \sim \mathcal{N}(\mathbf{0}, \mathbf{I})$ es ruido gaussiano muestreado de manera independiente.

\begin{importantbox}{Interpretación intuitiva de la fórmula de muestreo}
La actualización por paso de DDIM puede descomponerse en tres componentes:
\begin{enumerate}
    \item \textbf{Componente señal}: $\sqrt{\bar{\alpha}_{t-1}}\,\hat{\mathbf{x}}_0$——escala según la relación señal-ruido en el instante $t-1$, basándose en la mejor estimación actual de los datos limpios
    \item \textbf{Componente dirección}: vector que apunta desde $\hat{\mathbf{x}}_0$ hacia $\mathbf{x}_t$, conservando la "información direccional" del ruido
    \item \textbf{Componente aleatorio}: $\sigma_t \boldsymbol{\epsilon}_t$——ruido inyectado adicionalmente, cuya magnitud está controlada por $\sigma_t$
\end{enumerate}
Cuando $\sigma_t=0$, el componente aleatorio desaparece y todo el proceso se vuelve determinista.
\end{importantbox}

\subsection{Equivalencia de tres predictores}

El ponente señala que, al implementar el muestreo DDIM, se pueden utilizar tres parametrizaciones de red equivalentes:

\begin{enumerate}
    \item \textbf{Predictor de ruido} ($\boldsymbol{\epsilon}$-prediction): $\boldsymbol{\epsilon}_\theta(\mathbf{x}_t, t) \approx \boldsymbol{\epsilon}$, predice el ruido añadido
    \item \textbf{Predictor de datos limpios} ($\mathbf{x}_0$-prediction): $\hat{\mathbf{x}}_{0,\theta}(\mathbf{x}_t, t) \approx \mathbf{x}_0$, predice directamente los datos originales
    \item \textbf{Predictor de media} ($\boldsymbol{\mu}$-prediction): $\boldsymbol{\mu}_\theta(\mathbf{x}_t, t) \approx \boldsymbol{\mu}_{t-1}$, predice la media posterior
\end{enumerate}

\begin{knowledgebox}{Relación de conversión entre predictores}
Los tres predictores pueden convertirse unos en otros mediante transformaciones lineales simples:
$$
\hat{\mathbf{x}}_0 = \frac{\mathbf{x}_t - \sqrt{1-\bar{\alpha}_t}\,\boldsymbol{\epsilon}_\theta}{\sqrt{\bar{\alpha}_t}}, \quad
\boldsymbol{\epsilon}_\theta = \frac{\mathbf{x}_t - \sqrt{\bar{\alpha}_t}\,\hat{\mathbf{x}}_0}{\sqrt{1-\bar{\alpha}_t}}
$$
Independientemente de la parametrización utilizada en el interior de la red, el algoritmo de muestreo DDIM solo requiere una estimación de $\hat{\mathbf{x}}_0$. En la práctica, la parametrización $\boldsymbol{\epsilon}$-prediction es la forma más común (adoptada en el artículo original de DDPM), mientras que el $\mathbf{x}_0$-prediction también tiene aplicaciones en ciertos métodos mejorados.
\end{knowledgebox}

\subsection{Parametrización de $\sigma_t$ y casos especiales}

La elección de $\sigma_t$ controla el grado de aleatoriedad durante el proceso de muestreo. El artículo DDIM introduce un parámetro $\eta \geq 0$ para unificar la parametrización:

$$
\sigma_t(\eta) = \eta \cdot \sqrt{\frac{1-\bar{\alpha}_{t-1}}{1-\bar{\alpha}_t}} \cdot \sqrt{1 - \frac{\bar{\alpha}_t}{\bar{\alpha}_{t-1}}}
$$

\begin{table}[H]
\centering
\caption{Valores especiales del parámetro $\eta$ y el comportamiento de muestreo correspondiente}
\begin{tabular}{cll}
\toprule
$\eta$ & \textbf{Naturaleza del muestreo} & \textbf{Modelo equivalente} \\
\midrule
$0$ & Muestreo completamente determinista & DDIM (puramente determinista) \\
$1$ & Muestreo aleatorio, varianza consistente con DDPM & Equivalente a DDPM \\
$>1$ & Mayor aleatoriedad & Inyección de ruido más allá de DDPM \\
\bottomrule
\end{tabular}
\end{table}

\begin{warningbox}{DDPM es un caso especial de DDIM}
Cuando $\eta = 1$, la varianza de la distribución posterior $q_\sigma(\mathbf{x}_{t-1}|\mathbf{x}_t, \mathbf{x}_0)$ de DDIM es exactamente igual a la varianza posterior derivada en DDPM, $\tilde{\beta}_t = \frac{1-\bar{\alpha}_{t-1}}{1-\bar{\alpha}_t}\beta_t$. En este caso, DDIM se reduce al proceso de muestreo markoviano de DDPM. Sin embargo, dado que las ELBO de ambos son equivalentes, el uso de diferentes valores de $\eta$ para el muestreo no afecta el objetivo de entrenamiento—esto es clave para desacoplar el entrenamiento de la inferencia.
\end{warningbox}

\subsection{Resumen de la sección}

La fórmula de muestreo de DDIM se compone de tres términos: señal, dirección y aleatorio. El hiperparámetro $\sigma_t$ (o equivalentemente $\eta$) controla la magnitud de la aleatoriedad. $\eta=0$ proporciona muestreo determinista, mientras que $\eta=1$ degenera en DDPM. Independientemente de la elección de $\eta$, no es necesario volver a entrenar el modelo.
